% Осень26!A10; нумерация внутри столбца без пустых строк.
% Источник: лекция 3, §16 (строка 439)
\ticket{08}{Непрерывность вычисления обратного оператора}
\label{ticket:08}

\subsection*{\texorpdfstring{Следствия: \(\Omega_m\) открыто; непрерывность вычисления обратного оператора}{Следствия: \textbackslash Omega\_m открыто; непрерывность вычисления обратного оператора}}

\textbf{Следствие 1.} \(\Omega_m\) \textbf{открыто} в \(\mathrm{Lin}(\mathbb{R}^m,\mathbb{R}^m)\). Действительно (п. 1 теоремы §15): если \(L\in\Omega_m\), то весь шар радиуса \(1/\|L^{-1}\|\) вокруг \(L\) состоит из обратимых операторов --- «каждая точка внутренняя».

\textbf{Следствие 2.} Отображение \(L\mapsto L^{-1}\), вычисляющее обратный оператор, \textbf{непрерывно} на \(\Omega_m\) (задано только на \(\Omega_m\): «чтобы вычислять обратный, вы должны иметь обратимость; никуда вы его не расширите»).

\textbf{Доказательство (по Гейне).} Непрерывность в точке \(x_0\) по Гейне: \(x_n\to x_0\Rightarrow f(x_n)\to f(x_0)\). Пусть \(B_k\to L\) (в операторной норме), \(B_k\in\Omega_m\); нужно \(B_k^{-1}\to L^{-1}\). Так как \(\|L-B_k\|\to0\), при больших \(k\) имеем \(\|L-B_k\|<1/\|L^{-1}\|\), и по п. 3 теоремы §15: \[\|B_k^{-1}-L^{-1}\|\le\frac{\|L^{-1}\|}{\|L^{-1}\|^{-1}-\|L-B_k\|}\cdot\|L-B_k\|.\] Здесь \(L\) фиксировано («он не меняется»); при \(k\to\infty\) знаменатель дроби стремится к \(\|L^{-1}\|^{-1}>0\), значит дробь ограничена (стремится к \(\|L^{-1}\|^2\) --- «фактически вот эта штука в квадрате»), а множитель \(\|L-B_k\|\to0\). Итого правая часть \(\to0\), т.е. \(B_k^{-1}\to L^{-1}\).

«Прекрасно получился вывод: из сходимости \(B_k\) следует сходимость обратных --- это истолковывается как непрерывное взятие обратного отображения».
